\documentclass[11pt,a4paper]{article}
\usepackage[margin=1.8cm]{geometry}
\usepackage{amsmath,amssymb,tabularx,array,needspace,fontspec,enumitem,cancel}
\newcommand{\tb}{\cancel{||||}\,}
\setmainfont{FreeSerif}
\newfontfamily\urdufont[Script=Arabic]{FreeSerif}
\TeXXeTstate=1
\newcommand{\urdu}[1]{{\urdufont\beginR #1\endR}}
\setlength{\parindent}{0pt}
\setlength{\parskip}{0pt}
\renewcommand{\arraystretch}{1.15}
\pagestyle{plain}

% \mcq{number}{question}{A}{B}{C}{D}
\newcommand{\mcq}[6]{%
  \Needspace{4\baselineskip}%
  \par\medskip\noindent\textbf{#1.}\ #2\par\nopagebreak
  \noindent\begin{tabularx}{\linewidth}{@{\hspace{1.5em}}*{4}{>{\raggedright\arraybackslash}X}@{}}
    (A)~#3 & (B)~#4 & (C)~#5 & (D)~#6
  \end{tabularx}\par}

\begin{document}

\begin{center}
  {\Huge\bfseries The Hope Science Academy}\\[3pt]
  {\small Excellence is our standard not goal}\\[2pt]
  {\small 247-E-1, (Opposite LDA School) Johar Town, Lahore. PH: 0300-8194789, 0322-4157001}
\end{center}

\medskip
\noindent
\begin{tabularx}{\linewidth}{|X|X|X|X|}
\hline
\textbf{Student Name:} & \textbf{Roll Num:} & \textbf{Class:} 11th & \textbf{Paper Code:} \\ \hline
\textbf{Subject Name:} Statistics & \textbf{Time Allowed:} & \textbf{Total Marks:} 440 & \textbf{Exam Date:} 02-Nov-2026 \\ \hline
\multicolumn{4}{|l|}{\textbf{Exam Syllabus:} FULL BOOK} \\ \hline
\end{tabularx}

\bigskip
\noindent\textbf{Q1. Choose the correct answer.}\hfill\textbf{1$\times$440=440}

\mcq{1}{In the plural sense, statistics mean:}{Methods}{Numerical Data}{Sample values}{Population values}
\mcq{2}{In the singular sense, statistics means:}{Methodology}{Numerical data}{Sample values}{Count data}
\mcq{3}{Major area of statistics today is concerned with drawing conclusions from:}{Samples}{Populations}{Complete study}{Complete universe}
\mcq{4}{A collection of all the elements in a group is called:}{Population}{Sample}{Data}{Registration}
\mcq{5}{A collection some of the elements from a group is called:}{Census}{Population}{Sample}{Registration}
\mcq{6}{Statistic is characteristic calculated from:}{Population data}{Sample data}{Arrayed data}{Fictitious data}
\mcq{7}{The word statistics is at present used in:}{Two senses}{Three senses}{Four senses}{None of these}
\mcq{8}{If k is a constant then $\sum_{i=1}^{5}k$ :}{5k}{5 + k}{nk}{$k_{1}+k_{2}+k_{3}+k_{4}+k_{5}$}
\mcq{9}{Statistics came from the German word:}{Status}{Statista}{Statistik}{Statistique}
\mcq{10}{The mid point of group 5.5 ---7.5 is:}{6}{7}{7.5}{6.5}
\mcq{11}{In 60 --- 70, the lower limit is:}{50}{55}{60}{70}
\mcq{12}{A portion of population selected for study is called:}{Parameter}{Statistic}{Population}{Sample}
\mcq{13}{If "a" is a constant, then $\dfrac{5}{2}a$ is equal to:}{$a_{1}+a_{2}+a_{3}+a_{4}+a_{5}$}{a}{5a}{None}
\mcq{14}{Statistical laws are true:}{On the average}{Always}{In the long run}{A and C}
\mcq{15}{A numerical characteristics of a sample is called:}{Statistic}{Parameter}{Variable}{Sample}
\mcq{16}{The relative frequency multiplied by 100 is called:}{Percentage frequency}{Cumulative frequency}{Sample frequency}{Bivariate frequency}
\mcq{17}{A quantity calculated from a population is called:}{Frequency}{Statistics}{Parameter}{Sample}
\mcq{18}{Hourly temperature recorded by Weather Bureau represents:}{Discrete data}{Continuous data}{Qualitative data}{Secondary data}
\mcq{19}{Weight of whole earth is:}{Discrete variables}{Qualitative variable}{Continuous variable}{Constant}
\mcq{20}{Measurements usually provide:}{Discrete data}{Continuous data}{Qualitative data}{Primary data}
\mcq{21}{Counting or enumerations usually provide:}{Continuous data}{Qualitative data}{Discrete data}{Grouped data}
\mcq{22}{Number of accidents recorded yesterday in Lahore is a:}{Discrete variable}{Continuous variable}{Qualitative variable}{Constant}
\mcq{23}{The amount of milk given by a cow is a:}{Qualitative variable}{Discrete variable}{Continuous variable}{Constant}
\mcq{24}{Colour of hair is a:}{Continuous variable}{Discrete variable}{Qualitative variable}{Quantitative variable}
\mcq{25}{Smoking habits of residents of a city are:}{Qualitative data}{Quantitative data}{Discrete data}{Continuous data}
\mcq{26}{Continuous data are differentiated from discrete data in that:}{Discrete data classes are represented by fractions}{Continuous data classes may be represented by fractions}{Continuous data take on only whole numbers}{Discrete data can take on any real number}
\mcq{27}{Life of a T.V. tube is a:}{Discrete variable}{Continuous variable}{Qualitative variable}{Constant}
\mcq{28}{A constant can assume:}{Only one value}{More than one value}{Different values}{No value at all}
\mcq{29}{A statistic which is not measureable is called:}{A constant}{An attribute}{A variable}{A parameter}
\mcq{30}{The number 4.50001 rounded off to nearest unit (whole number) is:}{4}{5}{4.5}{4.01}
\mcq{31}{The number 5.56500 rounded off to nearest hundredth (two decimal places) is:}{5.57}{5.56}{6.00}{5.00}
\mcq{32}{The number 143.9500 rounded off to nearest tenth (one decimal place) is:}{143.9}{144.0}{143.0}{144}
\mcq{33}{The number 5.0435001 rounded off to nearest thousandth (three decimal places) is:}{5.043}{5.044}{5.050}{5.000}
\mcq{34}{The number 136.500 rounded off to nearest unit (whole number) is:}{136}{137}{136.5}{136.0}
\mcq{35}{Continuous variable can be measured at:}{Specific points}{All possible points}{No points}{Integer points}
\mcq{36}{A quantitative variable whose values are countable is called:}{Categorical variable}{Continuous variable}{Discrete variable}{None of these}
\mcq{37}{Weight of students in a class makes:}{Discrete data}{Continuous data}{Qualitative data}{Constant data}
\mcq{38}{The life time of fans, data is :}{Discrete}{Continuous}{Unchanged}{Qualitative}
\mcq{39}{Statistics must be:}{Comparable}{Not comparable}{Quantitative only}{None of these}
\mcq{40}{The phase of statistics that is concerned with the description and analysis of sample or population data is called:}{Inferential statistics}{Descriptive statistics}{Inductive statistics}{Sample statistics}
\mcq{41}{The phase of statistics that is concerned with the procedures and methodology or obtaining valid conclusions is called:}{Descriptive statistics}{Deductive statistics}{Inferential statistics}{Sample statistics}
\mcq{42}{Data used by an agency which originally collected them are:}{Primary data}{Raw data}{Secondary data}{Grouped data}
\mcq{43}{Registration is the source of:}{Primary data}{Secondary data}{Ogive}{Histogram}
\mcq{44}{Registration in the Population Census Reports are:}{Ungrouped data}{Secondary data}{Primary data}{Arrayed data}
\mcq{45}{Information recorded in its original form, whether counts or measurements, is referred to as:}{Continuous data}{Raw data}{Discrete data}{Arrayed data}
\mcq{46}{Issuing a national identity card is an example of:}{Census}{Registration}{Sampling}{Investigation through enumerators}
\mcq{47}{The data which have not undergone any statistical treatment are:}{Primary data}{Secondary data}{Discrete data}{Qualitative data}
\mcq{48}{The data which have already been collected by some one are called:}{Secondary data}{Primary data}{Arrayed data}{None of these}
\mcq{49}{How many sources are used for the collection of data:}{1}{2}{3}{4}
\mcq{50}{How many methods are used for the collection of data:}{1}{2}{3}{4}
\mcq{51}{Classification is the proces of arranging data according to:}{One characteristic}{Two or more characteristics}{Similar characteristics}{Different characteristics}
\mcq{52}{Important bases of classification are:}{Two}{Three}{Four}{More than four}
\mcq{53}{Data classified by attributes are called:}{Continuous data}{Quantitative data}{Qualitative data}{Grouped data}
\mcq{54}{Data which have been arranged in ascending or descending order are called:}{Raw data}{Grouped data}{Arrayed data}{Ungrouped data}
\mcq{55}{As a general rule, when arranging data statisticians tend to use:}{Less than six classes}{Between six and fifteen classes}{More than twenty classes}{Only fifteen classes}
\mcq{56}{Classification of data by their time of occurrence is called:}{Temporal or chronological classification}{Geographical classification}{Quantitative classification}{Qualitative classification}
\mcq{57}{Classification of data by attributes is called:}{Temporal classification}{Geographical classification}{Quantitative classification}{Qualitative classification}
\mcq{58}{Classification of data by quantitative characteristics is called:}{Qualitative classification}{Quantitative classification}{Geographical classification}{Temporal classification}
\mcq{59}{Classification of data according to locations or areas is called:}{Temporal classification}{Geographical classification}{Quantitative classification}{Qualitative classification}
\mcq{60}{Title of a table should be in:}{Lower case letters}{Capital letters}{Italic and lower case letters}{Twenty letters}
\mcq{61}{The process of arranging data into rows and c olumns is called:}{Frequency distribution}{Classification}{Tabulation}{Array}
\mcq{62}{A statistical table has at least:}{Five parts}{Four parts}{Three parts}{Two parts}
\mcq{63}{The section of table that contains the column caption is called:}{Box head}{Box plot}{Stub}{Body}
\mcq{64}{Two methods of data arrangement are:}{Array and frequency distribution}{Frequency distribution and histogram}{Array and frequency polygon}{Histogram and array}
\mcq{65}{Which of the following is an example of compressed data:}{Array}{Frequency distribution}{Histogram}{Ogive}
\mcq{66}{Give classes, 1 - 6, 6 - 10, ... class interval is:}{5.5}{3}{4}{5}
\mcq{67}{The grouped data is:}{Primary}{Secondary}{Raw data}{None of these}
\mcq{68}{An arrangement of data to show the frequency of occurrence is called:}{Frequency distribution}{Probability distribution}{Data array}{Cumulative distribution}
\mcq{69}{The average of lower and upper class limits is called:}{Class boundary}{Class frequency}{Class mark}{Class limit}
\mcq{70}{Frequency is denoted by:}{c}{q}{P}{f}
\mcq{71}{Mid point of the group 5.5 --- 7.5 is:}{6}{6.5}{7}{7.5}
\mcq{72}{When constructng a frequency distribution, the first step is:}{Divide the data into at least five classes}{Arrange the data into an array}{Decide on the type and number of classes for dividing the data}{None of these}
\mcq{73}{Which of the following statements is true of cumulative frequency polygons or ogives for a particular set of data:}{Both 'less than' and 'or more' curves have the same shape}{'Or more' curves slope up and to the right}{'Less than' curves slope down and to the right}{'Less than' curves slope up and to the right}
\mcq{74}{In constructing a frequency distribution for a sample, the number of classes depends on:}{The number of data points}{The range of the data collected}{The size of the population}{Both}
\mcq{75}{A relative frequency distribution presents frequencies in terms of:}{Fractions}{Whole numbers}{Percentages}{Both}
\mcq{76}{A graph of a cumulative frequency distribution is called:}{Histogram}{Frequency polygon}{Ogive}{None of these}
\mcq{77}{An ogive is a:}{Frequency curve}{Frequency polygon}{Cumulative frequency polygon}{Frequency histogram}
\mcq{78}{If a curve can be divided into two parts that are mirror images, it is called a:}{Skewed curve}{Symmetrical curve}{J-Shaped curve}{Frequency curve}
\mcq{79}{Mid points of top of the rectangular of histogram are joined to get:}{Polygon}{Frequency curve}{Ogive}{Histogram}
\mcq{80}{As the number of observations and classes increase, the shape of a frequency polygon:}{Tends to become increasingly smooth}{Tends to become jagged}{Stays the same}{Varies if data become more reliable}
\mcq{81}{Graphs of frequency distrubutions are used because:}{They have a long history in practical applications}{They attract attention to data pattern}{They account for biased or incomplete data}{None of these}
\mcq{82}{Continuous data are differentiated from discrete data in that:}{Discrete data classes are represented by fractions}{Continuous data classes may be represented by fractions}{Continuous data take on only whole numbers}{Discrete data can take on any real number}
\mcq{83}{If a curve has a longer tail to the right, it is called a:}{Positively skewed curve}{Negative skewed curve}{J-shaped curve}{Symmetrical curve}
\mcq{84}{If a curve has a longer tail to the left, it is called a:}{Symmetrical curve}{Positively skewed curve}{Negatively skewed curve}{None of these}
\mcq{85}{A chart in which total magnitude and its components are compared is called a:}{Component bar chart}{Pie chart}{Percentage component bar chart}{All of these}
\mcq{86}{Total angle of pie-chart is:}{270°}{300°}{320°}{360°}
\mcq{87}{Histogram is a graph of:}{Frequency distribution}{Time series}{Qualitative data}{Ogive}
\mcq{88}{The number of classes in a frequency distribution is obtained by dividing the range of variable by the:}{Total frequency}{Class interval}{Mid-points}{Relative frequency}
\mcq{89}{A frequency polygon is closed figure which is:}{One sided}{Two sided}{Three sided}{Many sided}
\mcq{90}{The frequency of a class divided by total frequency is called:}{Class frequency}{Cumulative}{Relative frequency}{Total frequency}
\mcq{91}{Brand of a soap is \_\_\_\_\_\_\_\_\_ variable:}{Quantitative}{Qualitative}{Imaginary}{Continuous}
\mcq{92}{The sum of rows or sum of columns of a bivariate frequency distribution are equal:}{Σf}{ΣX}{ΣfX}{100}
\mcq{93}{What is the major assumption we make when computing a mean from grouped date:}{All values are discrete}{Every value in a class is equal to the midpoint}{No value occurs more than once}{Each class contains exactly the same number of values}
\mcq{94}{When calculating the average rate of debt growth for a company, the correct mean to use is:}{Arithmetic mean}{Weighted arithmetic}{Geometric mean}{None of these}
\mcq{95}{Which of the following is the first step in calculating the median of a data set:}{Average the middle two values of the data set}{Array the data}{Determine the relative weights of the data values in terms of importance}{None of these}
\mcq{96}{Departure from symmetry is called:}{Kurtosis}{Skewness}{Dispersion}{None of these}
\mcq{97}{When a distribution is symmetrical and has one mode, the highest point on the curve is called the:}{Mode}{Median}{Mean}{All of these}
\mcq{98}{When referring to a curve that tails off to the left end, you would call it:}{Symmetrical}{Skewed to the right}{Positively skewed}{None of these}
\mcq{99}{In which of these cases would the mode be most useful as an indicator of central tendency:}{Every value in a data set occurs exactly once}{All but three values in a data set occur once; three values occur 5 times each}{All values n a data set occur 10 times each}{Every observation in a data set has the same value}
\mcq{100}{It is the reciprocal of the simple average of the reciprocal of all the values:}{A.M.}{G.M.}{H.M.}{Mode}
\mcq{101}{When referring to a curve whose longer tail is to the right, you would call it:}{Symmetrical}{Positively skewed}{Negatively skewed}{None of these}
\mcq{102}{If the mean is less than mode, the distribution is:}{Positively skewed}{Negatively skewed}{Symmetrical}{None of these}
\mcq{103}{The mean is affected by:}{Change of origin}{Scale of measurement}{Both}{None of these}
\mcq{104}{Sum of squares of deviations of the values is least when deviations are taken from:}{Median}{Mode}{Mean}{Harmonic mean}
\mcq{105}{When all the values in a series occur the same number of times, then it is not possible to compute the:}{Mean}{Median}{Mode}{Weighted mean}
\mcq{106}{The most central value of an arrayed data is:}{Mode}{Median}{Mean}{Harmonic mean}
\mcq{107}{The geometric mean of $X_{1}$ and $X_{2}$ is:}{$\dfrac{2}{X_{1}+X_{2}}$}{$\dfrac{2X_{1}X_{2}}{X_{1}+X_{2}}$}{$\dfrac{X_{1}+X_{2}}{2}$}{$\sqrt{X_{1}X_{2}}$}
\mcq{108}{Mode of 2, 10 and 7 is:}{10}{2}{7}{None of these}
\mcq{109}{The mean of the firstnnatural numbers is:}{n\{n + 1\}/2}{\{n + 1\}/2}{\{n - 1\}/2}{n/2}
\mcq{110}{The suitable average for the qualitative data is:}{Mean}{Mode}{Geometric mean}{Weighted mean}
\mcq{111}{If $\overline{X}=10$ and Y = 2X + 5, then $\overline{Y}=$ :}{20}{25}{30}{None of these}
\mcq{112}{If $\sum_{i=1}^{5}(X_{i}-20)=0$, then $\overline{X}=$ :}{0}{20}{5}{25}
\mcq{113}{Mode of the series 2, 2, 2, 3, 3, 3, 2, 3, 3, 4 is:}{3}{2 and 3}{4}{None of these}
\mcq{114}{Coding method is used for calculation of the:}{Median}{Mode}{Mean}{Weighted mean}
\mcq{115}{In a symmetrical distribution:}{$Q_{1}=Q_{3}$}{$P_{25}=P_{50}=P_{75}$}{A.M. = G.M. = H.M.}{A.M. = Med = Mode}
\mcq{116}{The mean of 10 observations is 10. All observations are increased by 10\%. The mean of the increased observations shall be:}{20}{11}{10}{100}
\mcq{117}{Which is appropriate average for finding the average speed of a journey:}{Mean}{Geometric mean}{Harmonic mean}{Weighted mean}
\mcq{118}{Which is the suitable average for calculating average percent increace in population:}{Median}{Geometric mean}{Harmonic mean}{Mean}
\mcq{119}{For a positively skewed distribution:}{Mean $>$ Mode}{Mode $>$ Mean}{Median $>$ mean}{None of these}
\mcq{120}{If any value in a series is zero, then we cannot calculate the:}{Mean}{Harmonic mean}{Mode}{Median}
\mcq{121}{If the values in a series are not of equal importance, we compute the:}{Median}{Harmonic mean}{Mean}{Weighted mean}
\mcq{122}{Which is the suitable average for calculating the average price at which articles are sold:}{Geometric mean}{Arithmetic mean}{Harmonic median}{Mode}
\mcq{123}{For a negatively skewed distribution:}{Mean $>$ Mode}{Mode $>$ Mean}{Mean $>$ Median}{None of these}
\mcq{124}{Which of the following measures cannot be calculated for the numbers 5, 8, 12, 6, 9, 13, 10:}{Median}{Mode}{Mean}{None of these}
\mcq{125}{The suitable average for shoe or collar size is:}{Geometric mean}{Arightmetic mean}{Mode}{Median}
\mcq{126}{Which pair of measures cannot be calculated when one of numbers in the series is zero:}{G.M. and A.M.}{G.M and H.M.}{H.M. and A.M.}{None of these}
\mcq{127}{The sum of deviations from arithmetic mean is $\Sigma(X-\overline{X})=$ :}{0}{1}{2}{3}
\mcq{128}{$\Sigma(X-20)=25$  $\Sigma(X-18)=0$ then arithmetic mean will be:}{25}{18}{38}{20}
\mcq{129}{If Mean = 40, Mode = 42, then distribution is:}{Positively skewed}{Negatively skewed}{Symmetrical}{All of these}
\mcq{130}{Which average cannot be computed if any value is less than zero:}{G.M.}{Median}{Mode}{A.M.}
\mcq{131}{Geometric mean of numbers 2, 4, 8 and 64 is:}{2}{4}{8}{64}
\mcq{132}{The Mean is based on:}{All the values}{Small values}{Extreme values}{Large values}
\mcq{133}{The arithmetic mean of 'n' values $y_{1}, y_{2}, \ldots, y_{n}$ is:}{$\dfrac{\Sigma fy}{\Sigma f}$}{$y_{1}+y_{2}+y_{3}+\ldots+y_{n}$}{$\dfrac{\Sigma y}{\Sigma yf}$}{$\dfrac{y_{1}+y_{2}+y_{3}+\ldots+}{n}$}
\mcq{134}{Median divides the data into:}{2 parts}{3 parts}{4 parts}{10 parts}
\mcq{135}{If a distribution has two modes, then it is called:}{Uni-model}{Bi-model}{Tri-model}{Multi model}
\mcq{136}{The most frequent value of the data if it exists is:}{A.M.}{G.M.}{Mode}{Median}
\mcq{137}{Mean is affected by the change of:}{Origin}{Scale}{Both a and b}{None}
\mcq{138}{If n=7 and $\Sigma(1/y)=1.929$, then H.M. is:}{5.77}{3.63}{5.89}{5.97}
\mcq{139}{If X=0, 2, 2, 4, 8, 10, then G.M. is:}{4}{8}{Zero}{10}
\mcq{140}{If $\Sigma(X-20)^{2}=0$, $\Sigma(X-18)^{2}=10$ then $\overline{X}$ could be:}{20}{18}{28}{25}
\mcq{141}{The median of -3, 0, -5, is:}{-3}{0}{-5}{Does not exist}
\mcq{142}{Which of the following average cannot be calculated from the observation 2, 2, 4, 4, 6, 6, 8, 8, 10, 10 :}{Mean}{Median}{Mode}{All of these}
\mcq{143}{In case of positively skewed distribution the extreme values lie in the:}{Middle}{Left tail}{Right tail}{Anywhere}
\mcq{144}{Which of the following averages is affected by extreme values:}{Arithmetic mean}{Median}{Mode}{All of these}
\mcq{145}{$\Sigma(Y-\overline{Y})=$ :}{0}{1}{Least}{> 0}
\mcq{146}{Which of the following is a relative measure of dispersion:}{Standard deviation}{Variance}{Coefficient of variation}{All of these}
\mcq{147}{If you divide the standard deviation of the values by the mean of the same values and multiply by 100, you will get:}{Variance}{Coefficient of variation}{Mean deviation}{Mean coefficient of dispersion}
\mcq{148}{The position square root of the variance of a distribution is the:}{Standard deviation}{Mean deviation}{Absolute deviation}{None of these}
\mcq{149}{The main advantage of using the range as a measure of dispersion is that:}{It is easy to calculate}{It is heavily influenced by extreme values}{It can change drastically from one data set to the next}{It is determined by only two points in the data set}
\mcq{150}{Why is it necessary to square the differences from the mean when computing the standard deviation:}{So that the extreme values will not affect the calculation}{Some of the differences will be positive and some will be negative}{The number of values involved is very small}{None of these}
\mcq{151}{The sum of absolute deviations is a minimum if these deviations are taken from the:}{Mean}{Median}{Mode}{All of these}
\mcq{152}{Range can be calculated in open-end classes:}{Never}{Always}{Often}{Seldom}
\mcq{153}{If s.d. (X) = 5, then s.d. $\left(\dfrac{2X+5}{2}\right)$ is equal to:}{5}{10}{15}{7.5}
\mcq{154}{The standard deviation is independent of:}{Change of origin}{Change of scale of measurement}{Change of origin and scale of measurement}{None of these}
\mcq{155}{If $\overline{X}=5$, which of the following expressions is a minimum:}{$\Sigma(X-25)^{2}$}{$\Sigma(X-5)^{2}$}{$\Sigma|X-5|$}{$\Sigma|X-25|$}
\mcq{156}{For moderately skewed distributions, approximately 68\% of the cases are included between:}{$\overline{X}-S$ and $\overline{X}+S$}{$\overline{X}-2S$ and $\overline{X}+2S$}{$\overline{X}-3S$ and $\overline{X}+3S$}{None of these}
\mcq{157}{For moderately skewed distributions, the approximate percentage of cases included between $\overline{X}-2S$ and $\overline{X}+2S$ is:}{99.7\%}{68\%}{95\%}{50\%}
\mcq{158}{Which of the following measures of dispersion is independent of the units employed:}{Standard deviation}{Quartile deviation}{Coefficient of variation}{Variance}
\mcq{159}{If X and Y are independent, then Var(X - Y) is equal to:}{Var \{X\} + Var \{Y\}}{Var \{X\} - Var \{Y\}}{Var \{X + Y\}}{Zero}
\mcq{160}{If X and Y are independent, then s.d. (X - Y) is equal to:}{$\sqrt{Var(X)-Var(Y)}$}{$\sqrt{Var(X)+Var(Y)}$}{S.D. \{X\} + S.D. \{Y\}}{S.D. \{X\}- S.D. \{Y\}}
\mcq{161}{For a moderately skewed distribution, which of the following empirical formula holds:}{$M.D.=\dfrac{4}{5}(S.D)$}{$Q.D.=\dfrac{2}{3}(S.D)$}{$Q.D.=\dfrac{5}{6}(M.D)$}{All of these}
\mcq{162}{The mean deviation is least if deviations are taken from:}{Median}{A.M.}{Mode}{G.M.}
\mcq{163}{If the third moment about mean is zero ($m_{3}=0$), then the distribution is:}{Mesokurtic}{Positively skewed}{Symmetrical}{Negatively skewed}
\mcq{164}{If the moment coefficient of kurtosis $b_{2}=3$, then the distribution is:}{Platykurtic}{Positively skewed}{Symmetrical}{Mesokurtic}
\mcq{165}{Which measure of dispersion is considered as the bes general purpose measure of dispersion:}{Range}{Semi- interquartile range}{Standard deviation}{Mean deviation}
\mcq{166}{To compare the variation of two or more than two series, we use:}{Mean absolute deviation}{Variance}{Coefficient of variation}{Corrected standard deviation}
\mcq{167}{The distribution is symmetrical if the moment coefficient of skewness $\sqrt{b_{1}}$ is:}{Negative}{Positive}{0}{3}
\mcq{168}{The distribution is mesokurtic if the moment coefficient of kurtosis $b_{2}$ is:}{Equal to 0}{Equal to 3}{Less than 3}{Greater than zero}
\mcq{169}{The distribution is positively skewed if:}{Mean $<$ Mode}{Mean $>$ Mode}{Mean $>$ Median}{Both}
\mcq{170}{Moment ratios $b_{1}$ and $b_{2}$ are:}{Expressed in original unit of the data}{Dimensionless quantities}{Independent of origin and scale of measurement}{Both}
\mcq{171}{For moderately skewed distributions, approximately 80\% of the cases are included between:}{$\overline{X}-S$ and $\overline{X}+S$}{$\overline{X}-2S$ and $\overline{X}+2S$}{$\overline{X}-3S$ and $\overline{X}+3S$}{None of these}
\mcq{172}{The sum of squares of deviations is a minimum if these deviations are taken from the:}{Median}{Mean}{Mode}{All of these}
\mcq{173}{Mean deviation is always:}{More than S.D.}{Equal to S.D.}{Less than S.D.}{None of these}
\mcq{174}{If X and Y are independent, then s.d. (X + Y) is equal to:}{$\sqrt{Var(X)+Var(Y)}$}{s.d.(X) + s.d.(Y)}{$\sqrt{Var(X).Var(Y)}$}{None of these}
\mcq{175}{Which of the following statements is correct:}{Every symmetrical curve is mesokurtic}{Standard deviation is the mean squared deviations from the mean}{The standard deviatoin of a constant is constant}{The second moment about zero equals variance}
\mcq{176}{The types of dispersion are:}{2}{3}{4}{5}
\mcq{177}{For a symmetrical distribution:}{$\beta_{1}=0$}{$\beta_{1}=3$}{$\beta_{2}=3$}{$\beta_{3}=3$}
\mcq{178}{Second moment about mean is:}{0}{1}{Variance}{Standard deviation}
\mcq{179}{The variance of 5, 5, 5, 5, 5 is:}{0}{25}{5}{125}
\mcq{180}{In symmetrical distribution if $Q_{1}=4$, $Q_{3}=12$ then median is:}{4}{6}{8}{Zero}
\mcq{181}{A disadvantage of range is that it is based on:}{Absolute deviation}{Square deviation}{Two extreme observation}{Upper and quartile}
\mcq{182}{Relative depression is of types:}{1}{4}{3}{2}
\mcq{183}{If Y = X + A, the range of Y= \_\_\_\_\_\_\_\_:}{Range}{Range}{Zero}{A}
\mcq{184}{For symmetrical distribution $\mu_{3}$ is:}{Zero}{1}{3}{4}
\mcq{185}{If the values of a variable are -2, -3, -5, -10 then range is:}{-12}{8}{-8}{0}
\mcq{186}{The measures of dispersion are changed by the change of:}{Origin}{Scale}{Both}{None of these}
\mcq{187}{M.D. of the values 4, 4, 4, 4 is:}{0}{4}{8}{12}
\mcq{188}{In Sheppard's correct $\mu_{2}$ is equal to:}{$\mu^{2}-\dfrac{h^{2}}{12}$}{$\mu^{2}+\dfrac{h^{2}}{12}$}{$\mu^{2}-\dfrac{h}{2}$}{$\mu+\dfrac{h}{12}$}
\mcq{189}{The mean deviation of dispersion can be negative:}{Never}{Always}{Sometimes}{Often}
\mcq{190}{For Leptokurtic distribution:}{$b_{2}>3$}{$b_{2}<3$}{$b_{2}=3$}{$b_{1}>3$}
\mcq{191}{First central moment is always:}{Zero}{1}{-1}{2}
\mcq{192}{The S.D. of 8, 8, 8, 8, 8 is:}{8}{$8^{2}$}{$\sqrt{8}$}{Zero}
\mcq{193}{Which set has the minimum variation:}{46, 48, 50}{30, 40, 50}{40, 50, 60}{48, 48, 49}
\mcq{194}{In a symmetrical distribution, the coefficient of skewness will always be:}{Negative}{Zero}{1}{-1}
\mcq{195}{If $Q_{3}=20$ and $Q_{1}=10$ the coefficient of quartile deviation is:}{3}{$\dfrac{1}{3}$}{$\dfrac{2}{3}$}{1}
\mcq{196}{The value of standard deviation changes by change of:}{Origin}{Scale}{Algebraic sign}{None}
\mcq{197}{The range of the scores 19, 3, 140, 25, 95 is:}{140}{3}{137}{143}
\mcq{198}{The range of the values -2, -4, -6 and -8 is:}{-6}{6}{-10}{-4}
\mcq{199}{Var (bX) = \_\_\_\_\_\_\_ :}{b Var}{$b^{2}$Var}{Var}{Var}
\mcq{200}{The variance expresses the variability of data in \_\_\_\_\_\_ as unit of data:}{Square of unit}{Square root of unit}{Same unit}{All of these}
\mcq{201}{If an index number calculation over 8 years with a base value of 100 gave an index for 1993 of 110, what would be the percentage relative for 1993:}{110}{90.0}{13.75}{880}
\mcq{202}{If all the values considered in calculating an index are of equal importance, the index is:}{Weighted}{Simple}{Unweighted}{None of these}
\mcq{203}{When the base year values are used as weights, the weighted average of relatives price index is the same as:}{The Paasche's index}{The Laspeyre's index}{The unweighted average of relatives price index}{None of these}
\mcq{204}{Simple index number involves commodities:}{One}{Two}{Four}{None of these}
\mcq{205}{If an unweighted average of relatives index $\dfrac{p_{n}}{p_{0}}\times100$ is calculated for each product in the composite, what is then done with these values to finish the calculation:}{The values are multiplied together}{The largest value is found}{The values are averaged}{None of these}
\mcq{206}{The weights used in a quantity index are:}{Percentages of total quantity}{Prices}{Average of quantities}{None of these}
\mcq{207}{The weights used in a price index are:}{Percentage of total price}{Quantities}{Average of prices}{None of these}
\mcq{208}{The base period can be described as a 'norma' period if:}{It is neither the peak nor the trough of a fluctuation}{It is the most recent period for which we have data}{It is the average of several consecutive periods}{None of these}
\mcq{209}{Link Relative is equal to:}{$\dfrac{P_{n}}{P_{0}}\times100$}{$\dfrac{P_{n}}{P_{n-1}}\times100$}{$\dfrac{P_{o}}{P_{n}}\times100$}{$\dfrac{P_{n}}{P_{n}}\times100$}
\mcq{210}{Which of the following methods uses quantities consumed in the current period when computing a weighted index:}{Laspeyre's method}{Paasche's method}{Marshall - Edge worth's method}{Fisher's method}
\mcq{211}{Which of the following methods uses quantities consumed in the base period when computing a weighted index:}{Laspeyre's method}{Paasche's method}{Fisher's method}{None of these}
\mcq{212}{Which of the following indices satisfies both the time reversal and factor reversal tests:}{Laspeyre's index}{Paasche's index}{Fisher's index}{Marshall - Edge worth}
\mcq{213}{Circular test is satisfied by:}{Laspeyre's index}{Paasche's index}{Fisher's index}{None of these}
\mcq{214}{To measure how much the cost of some variable changes over time, you would use:}{A value index}{An inflation index}{A quantity index}{None of these}
\mcq{215}{Commodities subject to considerable price variation should best be measured by:}{Quantity index}{Price index}{Value index}{None of these}
\mcq{216}{Theoretically best average used in the construction of composite index is:}{The arithmetic mean}{The geometric mean}{The median}{The harmomic mean}
\mcq{217}{Which of the following indices has an upward bias:}{Laspeyre's index}{Paasche's index}{Fisher's index}{None of these}
\mcq{218}{Paasche's Index Number:}{$\dfrac{\Sigma P_{n}q_{o}}{\Sigma P_{n}q_{o}}\times100$}{$\dfrac{\Sigma P_{n}q_{o}}{\Sigma P_{o}q_{n}}\times100$}{$\dfrac{\Sigma P_{n}q_{o}}{\Sigma P_{o}q_{o}}$}{$\dfrac{\Sigma P_{n}q_{n}}{\Sigma P_{o}q_{n}}\times100$}
\mcq{219}{Which is the most suitable average in chain base method:}{Arithmetic mean}{Median}{Geometric mean}{Weighted arithmetic mean}
\mcq{220}{The prices used in the construction of consumer price index numbers are:}{The retail prices}{The wholesale prices}{The fixed prices}{None of these}
\mcq{221}{The index given by $\dfrac{\Sigma p_{n}q_{n}}{\Sigma p_{o}q_{o}}\times100$ is:}{The Laspeyre's index}{The Paasche's index}{The value index}{None of these}
\mcq{222}{The consumer price index number is also called:}{The cost of living index numebr}{The retail price index number}{The wholesale price index number}{Both}
\mcq{223}{Which method of construction of consumer price index number is the Laspeyre's index number:}{Aggregate expenditure method}{Family budget method}{Both}{None of these}
\mcq{224}{Which of the following price indices are prepared by Federal Bureau of Statistics:}{Wholesale price index}{Consumer price index}{Sensitive price indicator}{All of these}
\mcq{225}{Which index number has a wide soope:}{Special}{General}{Price}{Quantity}
\mcq{226}{In chain base method, the base period is:}{Fixed}{Changed}{Constant}{None of these}
\mcq{227}{Base year weighted index number are:}{Laspeyre's}{Paasche's}{Fisher's}{C.P.I.}
\mcq{228}{Index for base period is always take as:}{100}{One}{200}{Zero}
\mcq{229}{Which formula is used in chain indices:}{$\dfrac{\Sigma P_{n}}{\Sigma P_{o}}\times100$}{$\dfrac{P_{n}}{\Sigma P_{n-1}}\times100$}{$\dfrac{P_{n}}{P_{o}}\times100$}{$\dfrac{\Sigma P_{n}q_{n}}{\Sigma P_{o}q_{o}}$}
\mcq{230}{Link relatives can be obtained dividing $P_{n}$ by:}{$P_{o}$}{$P_{n-1}$}{$q_{o}$}{$q_{n-1}$}
\mcq{231}{If Laspeyre's index number is 200, Paasche's index number is 200, the Fisher's index number is:}{100}{200}{Zero}{1000}
\mcq{232}{Price relatives is a ratio of current year price and :}{Base year quantity}{Previous year quantity}{Base year price}{Current year quantity}
\mcq{233}{An index that measures the change for a fixed time period is called:}{Chain base method}{Fixed base method}{Simple aggregative method}{Cost of living method}
\mcq{234}{How many basic types of index numbers:}{2}{3}{4}{5}
\mcq{235}{In fixed base method, the base period should be:}{Far away}{Normal}{Unreliable}{Abnormal}
\mcq{236}{In a fixed base method which period is taken always 100 :}{Preceding}{Following}{Base}{Current}
\mcq{237}{When all the commodities are not of equal importance, the index numbers are called:}{Simple}{Un-weighted}{Weighted}{Value}
\mcq{238}{$\dfrac{\Sigma p_{n}}{\Sigma p_{o}}\times100$ is called:}{Simple index}{Weighted index}{Simple aggregative index}{Price relative}
\mcq{239}{CPI falls in the category of :}{A simple index}{An aggregative}{An inflationary index}{Wholesale price index}
\mcq{240}{Which of the following formula is Laspeyre's price index number:}{$\dfrac{\Sigma p_{n}q_{0}}{\Sigma p_{0}q_{0}}\times100$}{$\dfrac{\Sigma p_{0}q_{n}}{\Sigma p_{0}q_{0}}\times100$}{$\dfrac{\Sigma p_{n}q_{n}}{\Sigma p_{0}q_{0}}\times100$}{All of these}
\mcq{241}{For computing chain index, we compute :}{Price relative}{Link relative}{Weighted indices}{None of these}
\mcq{242}{Consumer price index is obtained by:}{Paasche's formula}{Fisher's ideal formula}{Marshall Edgeworth formula}{Family Budget Method formula}
\mcq{243}{Composite index number involves commodities:}{One}{Two}{Three}{More than one}
\mcq{244}{An arrangement of all or some sets of objects in a definite order is called:}{Factorial}{Combinations}{Permutations}{None of these}
\mcq{245}{A set is usually denoted by a :}{Small letter}{Capital letter}{Both a and b}{None of these}
\mcq{246}{Two sets A and B are called overlapping sets if they have some :}{Different elements}{Common elements}{Negative element}{None of these}
\mcq{247}{The product set of A and B consists of all ordered pairs denoted by :}{$A\cdot B$}{$B\cdot A$}{$A\times B$}{None of these}
\mcq{248}{The complement of set A is :}{$U-A'$}{U - A}{$A'-U$}{None of these}
\mcq{249}{How many possible permutations can be formed form the wor ``PAKISTAN'':}{8!}{20160}{6!}{None of these}
\mcq{250}{An arrangement of all objects without carrying for the order is called :}{Permutation}{Binomial Expansion}{Combination}{None of these}
\mcq{251}{$C^{8}_{3}$ = \_\_\_\_\_\_\_\_ :}{20}{56}{6720}{None of these}
\mcq{252}{In the expansion of binomial, the coefficient of the first and the last term is always :}{2}{3}{1}{None of these}
\mcq{253}{The number of terms in the expansion of the binomial $(p+q)^{n}$ is:}{n}{n + 1}{n + 2}{None of these}
\mcq{254}{The difference between A and B shown by :}{A - B}{B - A}{A/B}{None of these}
\mcq{255}{$^{n}C_{r}$ = \_\_\_\_\_\_\_\_\_ :}{$C^{n-1}_{r}$}{$C^{n}_{n-r}$}{$C^{n}_{n-r-1}$}{None of these}
\mcq{256}{The product of n natural number is read as :}{$n^{n}$}{n!}{$\Sigma n$}{None of these}
\mcq{257}{$\dfrac{12!}{9!}$ = \_\_\_\_\_\_\_\_\_\_\_ :}{336}{1320}{3360}{None of these}
\mcq{258}{For $r>n$, $^{n}C_{r}$ = \_\_\_\_\_\_\_\_\_\_\_ :}{n}{1}{0}{not possible}
\mcq{259}{A set can be specified in two ways :}{Tabular and empty set form}{Tabular and set builder form}{Union and set builder form}{None of these}
\mcq{260}{A 300 years old person is an example of a :}{Difference set}{Null set}{Subset}{None of these}
\mcq{261}{B is a proper subset of A, if :}{A = B}{$A\subset B$}{$A\neq B$}{None of these}
\mcq{262}{In how many ways can 4 seats on a sofa set be occupied by 7 persons:}{60}{840}{28}{None of these}
\mcq{263}{The set of all subsets of a set is called :}{The power set}{Null set}{Product set}{None of these}
\mcq{264}{If $A\subset B$ and $B\subset C$, then A and B are :}{Difference}{Equal}{Nor}{None of these}
\mcq{265}{$C^{n}_{r}$ = \_\_\_\_\_\_\_\_\_\_\_ :}{$r!\times\,^{n}P_{r}$}{$\dfrac{1}{r!}\times P^{n}_{r}$}{$P^{n}_{r-1}$}{None of these}
\mcq{266}{The total number of subsets of a set containing n elements is :}{2n}{$2^{n}$}{$n^{2}$}{None of these}
\mcq{267}{A universal set is usually denoted by :}{A}{S}{U}{None of these}
\mcq{268}{$P^{8}_{6}$ = \_\_\_\_\_\_ :}{$\dfrac{8!}{(8-6)!}$}{$\dfrac{8!}{(8-6)!\,6!}$}{$\dfrac{8!}{6!}$}{None of these}
\mcq{269}{$^{n}C_{n-1}$ = \_\_\_\_\_\_\_\_\_ :}{n}{r}{n}{None of these}
\mcq{270}{The permutation of the word ``ARRANGEMENT'' is :}{2449800}{248900}{2494800}{None of these}
\mcq{271}{The permutation of the word ``PAKPATTAN'' is :}{11520}{15120}{15210}{None of these}
\mcq{272}{The permutation of the word ``STATISTICS'' is :}{500000}{50400}{54000}{None of these}
\mcq{273}{Four books can be arranged in a self in :}{8 ways}{12 ways}{16 ways}{24 ways}
\mcq{274}{If one event is unaffected by the outcome of another event, the two events are said to be :}{Dependent}{Independent}{Mutually exclusive}{Both}
\mcq{275}{Symbolically, a marginal probability is :}{P\{AB\}}{P\{A\}}{$P(A\,|\,B)$}{Both a and b}
\mcq{276}{The simple probability of occurrence of an event is called the :}{Joint probability}{Conditional probabiliy}{Marginal probability}{Subjective probability}
\mcq{277}{Why are the outcomes of a coin tossing mutually exclusive:}{The outcome of any toss is not affected by the outcome of those preceding it}{Both a head and a tail cannot turn up on any one toss}{The probability of getting a head and the probability of getting a tail is the same}{All of these}
\mcq{278}{What is the probability that a value chosen at random from a particular population is larger than the median of the population:}{0.25}{0.5}{1.0}{0.67}
\mcq{279}{What is the probability that a ball drawn at random from the bag is :}{0.1}{0.4}{1.0}{Cannot be determined from given information}
\mcq{280}{When two dice are rolled, the number of possible sample points is :}{6}{12}{36}{48}
\mcq{281}{The probability of drawing a king of spade from a pack of 52 cards is :}{$\dfrac{1}{4}$}{$\dfrac{1}{13}$}{$\dfrac{1}{26}$}{$\dfrac{1}{52}$}
\mcq{282}{If P(A or B) = P(A), then :}{+ P}{is the joint probability of A and B}{P}{None of these}
\mcq{283}{If a Venn diagram is drawn for events A and B which are mutually exclusive, which of the following would always be true of A and B:}{Their parts of the rectangle will overlap}{Their parts of the rectangle will be equal in area}{Their parts of the rectangle will not overlap}{None of these}
\mcq{284}{When two coins are tossed simultaneously, P(one head) is :}{$\dfrac{1}{4}$}{$\dfrac{1}{2}$}{$\dfrac{3}{4}$}{1.0}
\mcq{285}{When three coins are tossed simultaneously, P(3 heads) is :}{$\dfrac{3}{8}$}{$\dfrac{1}{2}$}{$\dfrac{1}{4}$}{$\dfrac{1}{8}$}
\mcq{286}{The probability of drawing two aces from a pack of 52 cards with replacement is :}{$\dfrac{1}{4}$}{$\dfrac{1}{16}$}{$\dfrac{1}{169}$}{$\dfrac{1}{256}$}
\mcq{287}{For two mutually exclusive events A and B, P(A) = 0.2 and P(B) = 0.4, then $P(A\cup B)$ is :}{0.8}{0.2}{0.6}{0.5}
\mcq{288}{For two independent events A and B, P(A) = 0.6 and P(B) = 0.3, then $P(A\cap B)$ is :}{0.9}{0.12}{0.5}{0.18}
\mcq{289}{If A and B are mutually exclusive events, then $P(A\cup B)=P(A)+P(B)$. How does the calculation of $P(A\cup B)$ change if A and B are not mutually exclusive:}{+ P}{] must be divided by $P(A\cap B)$}{[P}{[P}
\mcq{290}{Assume that a single fair dice is rolled once. Which of the following is true:}{The probability of rolling a number greater than '1' is 1 - P\{'1' is rolled\}}{The probability of rolling a '3' is 1 - P\{1, 2, 4, 5 or 6 is rolled\}}{The probability of rolling a '5 or 6' is higher than the probability of rolling a '3 or 4'}{Both a and b}
\mcq{291}{If $P(B|A)=0.25$ and $P(A\cap B)=0.20$, then P(A) is :}{0.05}{0.8}{0.95}{0.75}
\mcq{292}{The probability of red card out of 52 cards is :}{$\dfrac{1}{4}$}{$\dfrac{1}{2}$}{$\dfrac{4}{52}$}{Zero}
\mcq{293}{When two coins are tossed simultaneously, the probability of at most one head is:}{$\dfrac{1}{4}$}{$\dfrac{1}{2}$}{$\dfrac{3}{4}$}{None of these}
\mcq{294}{A letter is chosen at random the word STATISTICS. The probability of getting a vowel is :}{$\dfrac{1}{5}$}{$\dfrac{3}{10}$}{$\dfrac{1}{2}$}{$\dfrac{2}{5}$}
\mcq{295}{The probability of getting one red ball from a bag containing 4 red, 3 white and 3 black balls is :}{$\dfrac{3}{10}$}{$\dfrac{1}{5}$}{$\dfrac{2}{5}$}{$\dfrac{1}{2}$}
\mcq{296}{$P(A\cap B)=P(A|B)P(B)$, then A and B are :}{Independent}{Mutually exclusive}{Dependent}{None of these}
\mcq{297}{The probability of getting two red balls with replacement from a bag containing 4 red, 3 white and 3 black balls is :}{$\dfrac{4}{25}$}{$\dfrac{1}{25}$}{$\dfrac{9}{100}$}{$\dfrac{2}{25}$}
\mcq{298}{Ten numbered balls are placed in an urn. Numbers 1 - 4 are red and numbers 5 - 10 are blue. The probability that a ball drawn at random from the urn is blue is:}{0.1}{0.4}{0.6}{1.0}
\mcq{299}{For two independent events A and B, P(A) = 0.6 and P(B) = 0.3. Then $P(A\cup B)$ is:}{0.9}{0.18}{0.72}{None of these}
\mcq{300}{Which of the following approaches to probability would you apply to solve the problems given afterwards:}{Classical approach}{Relative Frequency approach}{Subjective approach}{Axiomatic approach}
\mcq{301}{The probability of getting an odd number when a balanced dice is rolled is :}{$\dfrac{1}{6}$}{$\dfrac{1}{4}$}{$\dfrac{1}{2}$}{$\dfrac{1}{3}$}
\mcq{302}{If $P(B|A)=P(B)$, then A and B are :}{Mutually exclusive}{Dependent}{Independent}{None of these}
\mcq{303}{If $P(A\cup B)=P(A)+P(B)$, then A and B are :}{Mutually exclusive}{Independent}{Not mutually exclusive}{None of these}
\mcq{304}{The probability of drawing any one spade card is :}{$\dfrac{1}{32}$}{$\dfrac{1}{13}$}{$\dfrac{4}{13}$}{$\dfrac{1}{4}$}
\mcq{305}{In the tossing of two perfect coins the probability at least one head occur is :}{$\dfrac{1}{2}$}{$\dfrac{1}{4}$}{$\dfrac{3}{4}$}{1}
\mcq{306}{The probability of an event cannot be :}{=0}{$>$0}{=1}{$<$0}
\mcq{307}{If two coins are tossed, the probability of getting one head and one tail is :}{$\dfrac{1}{4}$}{$\dfrac{2}{4}$}{3/4}{$\dfrac{2}{3}$}
\mcq{308}{A measure of the change that an uncertain event will occur :}{An experiment}{An event}{A probability}{A trial}
\mcq{309}{When a dice and a coin are rolled together, all possible outcomes are :}{6}{2}{36}{12}
\mcq{310}{Two events A and B are called mutually exclusive if :}{$A\cup B=\phi$}{$A\cap B=\varphi$}{$A\cap B=S$}{$A\cap B=1$}
\mcq{311}{A set of numerical values assigned to a sample space is called :}{Random sample}{Random variable}{Random numbers}{Random experiment}
\mcq{312}{Events with equal probabilities are called :}{Mutually exclusive events}{Exhaustive events}{Equally likely events}{Simple events}
\mcq{313}{If A and B are mutually exclusive events then $P(A\cap B)$ is :}{Zero}{1}{0.5}{0.8}
\mcq{314}{The collection of all possible outcome of a random experiment is called :}{Sample point}{Sure event}{Sample space}{Simple event}
\mcq{315}{A student solved 25 questions from first 50 questions of a book to be solved. The prob. that he will solve the remaining all question :}{0.25}{0.5}{1}{0}
\mcq{316}{If $P(A)=\dfrac{1}{3}$, $P(B)=\dfrac{2}{3}$ if A and B are independent then $P(A\cap B)$ is :}{$\dfrac{2}{6}$}{2/9}{0}{$\dfrac{1}{9}$}
\mcq{317}{The result of no interest of an experment is called :}{Constant}{Event}{Failure}{Success}
\mcq{318}{Which is the impossible event when a dice is rolled:}{5 or 6}{6 or 7}{2 or 3}{1}
\mcq{319}{When two dice are rolled, the maximum total on the two faces of dice will be :}{6}{36}{12}{2}
\mcq{320}{The probability of an impossible event is :}{Positive}{Zero}{Negative}{1}
\mcq{321}{Two coins are tossed, the probability of two tails is equal to :}{$\dfrac{1}{4}$}{$\dfrac{1}{2}$}{$\dfrac{3}{4}$}{1}
\mcq{322}{The range of probability is between :}{0 to 2}{-1 to +1}{0 to $\infty$}{0 to 1}
\mcq{323}{If $\Sigma(XY)=\Sigma(X)(Y)$, then X and Y are :}{Dependent}{Independent}{Biased}{Unbiased}
\mcq{324}{Tossing two dice possible samples are :}{36}{12}{6}{2}
\mcq{325}{The term ``event'' is used for :}{Time}{Subset of the sample space}{Probability}{Total number of outcomes}
\mcq{326}{Which of the following cannot be probability of an event:}{0}{1}{0.32}{1.05}
\mcq{327}{When a pair of dice is rolled, the sum of uppermost dots vary from :}{0 to 10}{1 to 11}{2 to 19}{2 to 12}
\mcq{328}{Every random experiment result in outcomes:}{two or more}{three or more}{four or more}{five or more}
\mcq{329}{When a pair of dice is thrown, the sum of upturned dots vary from:}{0 to 10}{2 to 12}{1 to 11}{3 to 13}
\mcq{330}{Random variable are donated by :}{L, M or N}{A,B or C}{X,Y or Z}{P,Q or R}
\mcq{331}{The values of random variables are donated by:}{p,q or r}{l,m or n}{a,b or c}{x,y or z}
\mcq{332}{In pair of dice, there are outcomes:}{32}{36}{38}{48}
\mcq{333}{Each digit \{0, 1, 2, .....9\} has probability:}{0.4}{0.2}{0.1}{0.3}
\mcq{334}{\{0,1,2......9\} are also called:}{random variables}{random outcomes}{random experiments}{random digits}
\mcq{335}{We can get random numbers from:}{games of change}{readin books}{different people}{reading tables}
\mcq{336}{Tables of random numbers are widely :}{acceptable}{available}{readable}{countable}
\mcq{337}{In pseudo random techniques, the digits will re-appear in the:}{cycle}{manner}{order}{raoutine}
\mcq{338}{Psedudo random digits are not truely:}{constructed}{possible}{associated}{random}
\mcq{339}{Pseudo random digits are completely:}{Probabillistic}{realistic}{deterministic}{none of these}
\mcq{340}{$ax_{n}+b$ (mod m) is equal to:}{$x_{n+1}$}{$x_{n-1}$}{$x_{n}$}{$x_{n+2}$}
\mcq{341}{The seed $x_{o}$ (a starting number is) :}{A whole number}{A natural number}{An integral number}{A rational number}
\mcq{342}{The range of psedou random digit is from:}{0 to m +1}{a natural number}{An integer}{a ratiomal number}
\mcq{343}{Random numbers are widely used in:}{analyzation}{simulation number}{criticization}{experimentation}
\mcq{344}{Probability proportional to size is abbreviated as:}{PPS}{SPP}{PPP}{PSP}
\mcq{345}{The random number o to 9 ae called random numbers:}{1-digit}{2-digit}{3-digit}{4-digit}
\mcq{346}{PPS is used when we have information about:}{random variables}{probabillities}{probability distribution}{index numbers}
\mcq{347}{00-999 are called random numbers:}{1-digit}{2-digit}{-digit}{4-digit}
\mcq{348}{000-999 are called random numbers:}{1-digit}{2-digit}{3-digit}{4-digit}
\mcq{349}{If the set of values comprises isolated points on a real line, then the variable is called:}{discrete}{continuous}{given}{random}
\mcq{350}{A random variable may be discrete or:}{experimental}{functional}{given}{continuous}
\mcq{351}{A random variable may be discrete or:}{experimental}{functional}{given}{continuous}
\mcq{352}{Random variable is also called:}{discrete variabel}{continuous variable}{chance variable}{vaiable}
\mcq{353}{If X and Y are random variable then E(x+y) = E(x) + E(y) is equal to:}{E(x) +E(y)}{E(x) -E(y)}{X-E(y)}{E(x)-Y}
\mcq{354}{A balanced diet is rolled, the probabillity of getting odd number is:}{1/2}{1/4}{1/6}{1/3}
\mcq{355}{If two dice ar rolled, the no. of the possible outcomes are:}{6}{36}{1}{difficult to all}
\mcq{356}{A ....... random variable can assume each and every value within an interval:}{continuous}{qualitatve}{drscrete}{none of these}
\mcq{357}{A variable that can assume only some specific value within a given range is called:}{continuous r.v}{discrete r.v}{sample space}{difficult to tell}
\mcq{358}{Random numbers are squence of digits from the set:}{\{0,,2,........,9\}}{\{1,2,3,......9\}}{\{1,2,3,......10\}}{\{2,4,6,......10\}}
\mcq{359}{A random variable which assume finite number of value is called:}{continuous random variable}{deviated variable}{discrete random variable}{all of these}
\mcq{360}{A tabular arrangment of probabillities against each possible value of discrete random variable is known as:}{distribution function}{density function}{probability density function}{none of these}
\mcq{361}{A variable that is the result of a measuring process is called a:}{uniform variable}{discrete variable}{continuous variable}{all of these}
\mcq{362}{Every random experiment has:}{only one outcome}{at least 2 outcomes}{0}{none of these}
\mcq{363}{A discrete random variable takes only:}{infinte values}{countable values}{uncountable values}{none of these}
\mcq{364}{The total are under the density curve is:}{0}{1}{2}{none of these}
\mcq{365}{If x is a countinuous random variable, then the peobillity that x take a specific value is:}{1}{0}{$2^{2}$}{none of these}
\mcq{366}{A random variable may not assume:}{positive values}{nagative values}{zero}{none of these}
\mcq{367}{A discrete-probabillity function 'f(x)' is always:}{non-negative}{negative}{0}{none of these}
\mcq{368}{A random variable may be disccret or:}{experimental}{functional}{given}{continuous}
\mcq{369}{A discrete probability distribution may be represented by :}{A table}{A graph}{A mathematical equation}{All of these}
\mcq{370}{A continuous probability distribution may be represented by :}{A table}{A graph}{A mathematical equation}{Both}
\mcq{371}{If the random variable X denotes the number of heads when three distinct coins are tossed, then X assumes the values :}{0, 1, 2, 3,}{1, 3, 3, 1}{1, 2, 3}{None of these}
\mcq{372}{If X is a discrete random variable, then the function f(x) is :}{A probability function}{A probability density function}{A density function}{A distribution function}
\mcq{373}{If X is a continuous random variable, then the function f(x) is :}{A probability function}{A probability density function}{A density function}{Both}
\mcq{374}{The distribution function F(x) is equal to :}{P}{$P(X\leq x)$}{$P(X\geq x)$}{None of these}
\mcq{375}{The expected value of a discrete random variable is :}{Always an integer}{Always one of the values that the random variable can assume}{An interval of values}{None of these}
\mcq{376}{If X and Y are random variables, then E(X - Y) is equal to :}{E\{X\} + E\{Y\}}{E\{X\} - E\{Y\}}{X - E\{Y\}}{E\{X\} - Y}
\mcq{377}{If X and Y are independent random variables, then Var (X - Y) is equal to :}{Var \{X\} - Var \{Y\}}{Var \{X\} + Var \{Y\}}{$E\{X-Y\}^{2}$}{$E\{X+Y\}^{2}$}
\mcq{378}{If X and Y are independent random variables, then S.D. (X - Y) is equal to :}{$\sqrt{Var(X)+Var(Y)}$}{$\sqrt{Var(X)-Var(Y)}$}{S.D.}{S.D.}
\mcq{379}{If X and Y are independent random variables, then E(XY) is equal to :}{E\{XY\}}{E\{X\} E\{Y\}}{XE\{Y\}}{YE\{X\}}
\mcq{380}{If X is a random variable and a and b are constants, then Var (a + bX) is equal to:}{Var \{a\} + Var \{bX\}}{Var \{a\} + b Var \{X\}}{$b^{2}$Var \{X\}}{b Var \{X\}}
\mcq{381}{If X is a random variable and a and b are constants, then S.D. (a - bX) is equal to:}{S.D.\{a\} - S.D. \{bX\}}{S.D. \{a\} - b S.D. \{X\}}{$b\sqrt{Var(X)}$}{None of these}
\mcq{382}{If X and Y are random variables and a and b are constants, the E(aX + bY) is equal to :}{E\{aX\} + E\{bX\}}{aE\{X\} + bE\{Y\}}{\{a + b\}\{E\{X\} + E\{Y\}\}}{None of these}
\mcq{383}{If E(X) = 4, then arithmetic mean will be :}{4}{8}{0}{1}
\mcq{384}{If 'C' is a non-random variable, then E(C) is :}{c}{0}{1}{x}
\mcq{385}{Given Var(X) = 3 and Var(Y) = 5. If X and Y are independent variables then Var(X - Y)= :}{-2}{8}{24}{2}
\mcq{386}{$E(X-\mu)^{2}$ = \_\_\_\_\_\_\_\_ :}{Zero}{S.D.}{Variance}{Mean deviation}
\mcq{387}{An expected value of a random variable is equal to its :}{Mean}{Variance}{S.D.}{Co-variance}
\mcq{388}{If y=5x+10, then S.D.(y) is :}{5 S.D.\{x\}}{25 S.D. \{x\}}{10 S.D.\{x\}}{5 S.D. \{x\} + S.D. \{10\}}
\mcq{389}{If X is a random variable, then E(aX+b) is equal to :}{E\{X\}}{aE\{X\}}{aE\{X\} + b}{E\{X\} + a}
\mcq{390}{If $Y_{i}=aX_{i}+b$, where 'a' and 'b' are any two constants the Var(y)= :}{a Var\{X\} + b}{Var\{X\} + b}{$a^{2}$ Var\{X\}}{a Var\{X\}}
\mcq{391}{For a random variable 'X' $\Sigma P(x)$ is equal to :}{Zero}{1}{Greater than 1}{Less than 1}
\mcq{392}{The probability that a continuous random variable 'X' takes on specific value of x is :}{Greater than zero}{Less than zero}{Equal to zero}{0 to 1}
\mcq{393}{Probability density function is the probability function of \_\_\_\_\_ random variable :}{Continuous}{Discrete}{Qualitative}{None}
\mcq{394}{The probability density function P(x) cannot exceed :}{Zero}{One}{Mean}{Infinity}
\mcq{395}{If $f(x)=\dfrac{1}{10}$, X = 10 then E(X) is :}{Zero}{6/8}{1}{-1}
\mcq{396}{Which of the following is not possible in a probability distribution :}{f}{$\Sigma f(x)=1$}{$\Sigma xf(x)=2$}{f}
\mcq{397}{Var(x) is equal to :}{$[E(x)]^{2}-E(x^{2})$}{$E(x^{2})-[E(x)]^{2}$}{$E(x^{2})$}{E\{x\}}
\mcq{398}{In a discrete probability distribution the sum of all the probabilities is always:}{0}{1}{-1}{∞}
\mcq{399}{For a given binomial distribution withn fixed, if p $<$ 0.5, then:}{The binomial distribution will be skewed to the left}{The binomial distribution will be skewed to the right}{The binomial distribution will be symmetric}{None of these}
\mcq{400}{If the probability of success p = 0.4 for a particular Bernoulli trial, the expression $\dfrac{7!}{3!\,4!}(0.4)^{3}(0.6)^{4}$ gives the probability of getting:}{Exactly three successes in seven trials}{Exactly four successes in seven trials}{Three or more successes in seven trials}{Four or more successes in seven trials}
\mcq{401}{The mean of a binomial distribution depends on :}{Probability of success}{Probability of failure}{Number of trials}{Both a and c}
\mcq{402}{The mean of a binomial distribution depends on :}{Probability of success}{Probability of failure}{Number of trials}{Both 2}
\mcq{403}{The standard deviation of a binomial distribution depends on :}{Probability of success}{Probability of failure}{Number of trials}{Both}
\mcq{404}{Which of the following can never be described by a binomial distribution:}{The number of defective items produced by an assembly process}{The amount of water used by a single household}{The number of students in the class who can answer this question}{All of these can always be described by a binomial distribution}
\mcq{405}{Which of the following is true for the binomial distribution b(x; n, p):}{mean > variance}{mean < variance}{mean = variance}{mean = standard deviation}
\mcq{406}{The mean, median and mode of the binomial distribution b(x; n, p) will be equal when :}{p = 0.5}{p < 0.5}{p > 0.5}{None of these}
\mcq{407}{The parameters of the hypergeometric distribution h(x; N, n, k) are :}{x, N, k}{N, n, k}{x, n, k}{x, N, n, k}
\mcq{408}{The number of possible outcomes in a Bernoulli trial is :}{Three}{Four}{Two}{One}
\mcq{409}{A binomial random variable is a(an) :}{Continuous random variable}{Discrete random variable}{Independent variable}{Dependent variable}
\mcq{410}{A hypergeometric random variable is a (an) :}{Independent variable}{Continuous random variable}{Discrete random variable}{None of these}
\mcq{411}{In which distribution the probability of success remains constant from trial to trial:}{Hypergeometric distribution}{Binomial distribution}{Sampling distribution}{Continuous distribution}
\mcq{412}{In which distribution the probability of success changes from trial to trial:}{Binomial distribution}{Sampling distribution}{Hypergeometric distribution}{Continuous distribution}
\mcq{413}{In which distribution the successive trials are with replacement:}{Hypergeometric distribution}{Binomial distribution}{Continuous distribution}{None of these}
\mcq{414}{In which distribution the successive trials are without replacement:}{Hypergeometric distribution}{Binomial distribution}{Continuous distribution}{None of these}
\mcq{415}{Both binomial and hypergeometric distributions are :}{Continuous probability distribution}{Discrete probability distributions}{Neither continuous nor discrete probability distributions}{Bivariate distributions}
\mcq{416}{A fair coin is tossed four times. The probability of getting four heads is :}{$\dfrac{1}{4}$}{$\dfrac{1}{2}$}{$\dfrac{1}{16}$}{1}
\mcq{417}{A fair die is rolled three times. The probability of getting three 'aces' is :}{$\dfrac{1}{3}$}{$\dfrac{1}{6}$}{$\dfrac{1}{216}$}{$\dfrac{1}{27}$}
\mcq{418}{A fair coin is tossed five times. The probability of getting zero head is :}{0}{$\dfrac{1}{32}$}{$\dfrac{1}{5}$}{$\dfrac{1}{2}$}
\mcq{419}{A fair coin is tossed five times. The probability of getting zero head is :}{0}{$\dfrac{1}{32}$}{$\dfrac{1}{5}$}{$\dfrac{1}{2}$}
\mcq{420}{In hypergeometric distribution, the trials are :}{Independent}{Dependent}{Independent and dependent}{None of these}
\mcq{421}{In a binomial experiment, the successive trials are :}{Dependent}{Independent}{Mutually exclusive}{Fixed}
\mcq{422}{The binomial probability distribution is symmetrical when :}{P = q}{P $<$ q}{P $>$ q}{np $>$ npq}
\mcq{423}{The percentage of observations lying within the limits $\overline{X}\pm3S$ in the normal distribution is :}{68.26\%}{95.44\%}{70.00\%}{99.73\%}
\mcq{424}{In a hypergeometric distribution N = 6, n = 2, K = 3 then mean = :}{2}{3}{1}{4}
\mcq{425}{In binomial distribution the range of X is :}{Zero to $\infty$}{Zero ton}{-1 to $\infty$}{Undefined}
\mcq{426}{In hyper-geometric distributionn is :}{Changed}{Zero}{Fixed}{Variable}
\mcq{427}{In a binomial, n=20, p=3/5, then variance of this distribution is :}{60}{12}{4.8}{0}
\mcq{428}{If n = 10 and p = 0.6 the $P(X\geq0)$ is :}{0.5}{0.6}{1.0}{1.2}
\mcq{429}{Binomial distribution is positive skewed when :}{P = q}{P > q}{P < q}{$P=\dfrac{1}{2}$}
\mcq{430}{If $P=q=\dfrac{1}{2}$, then distribution is called :}{Symmetrical}{Positively}{Skewed}{Negatively}
\mcq{431}{The hyper-geometric distribution has \_\_\_\_\_\_ parameters :}{2}{3}{1}{4}
\mcq{432}{The probability of failure is equal to :}{P}{1 - q}{P - 1}{1 - P}
\mcq{433}{The mean of the hyper geometric distribution is :}{$\dfrac{Nn}{k}$}{$\dfrac{Nk}{k}$}{$\dfrac{nk}{N}$}{$\dfrac{n+k}{N}$}
\mcq{434}{The mean and variance of the binomial distribution are :}{np and $\sqrt{npq}$}{np and npq}{np and nq}{n and p}
\mcq{435}{In a binomial experiment with three trials, the variable can take :}{2 values}{3 values}{4 values}{5 values}
\mcq{436}{The number of trial in binomial distribution is :}{Not fixed}{Fixed}{Large}{Small}
\mcq{437}{In a hypergeometric distribution :}{Mean $>$ variance}{Mean $<$ variance}{Mean = variance}{Mean = zero}
\mcq{438}{The binomial distribution is negatively skewed if :}{$p<\dfrac{1}{2}$}{$p=\dfrac{1}{2}$}{$p>\dfrac{1}{2}$}{p = 1}
\mcq{439}{If N = 40, n = 5, k = 4 then mean of hypergeometric distribution is :}{1}{1/2}{1/4}{1/3}
\mcq{440}{Binomial distribution is negatively skewed if :}{p $<$ q}{p $>$ q}{p = q}{np = npq}

\end{document}
